\documentclass[11pt]{article}
\usepackage[margin=1in]{geometry}
\usepackage{amsmath,amssymb,amsthm,booktabs}
\usepackage[hidelinks]{hyperref}
\setlength{\emergencystretch}{2em}
\newtheorem{theorem}{Theorem}
\newtheorem{lemma}[theorem]{Lemma}
\newtheorem{corollary}[theorem]{Corollary}
\newcommand{\mr}{\operatorname{mr}}
\newcommand{\rank}{\operatorname{rank}}
\newcommand{\tr}{\operatorname{tr}}
\newcommand{\adj}{\operatorname{adj}}
\newcommand{\C}{\mathbb C}
\newcommand{\Q}{\mathbb Q}
\newcommand{\Z}{\mathbb Z}
\title{Reflective augmentation rank below meridional rank}
\author{Internal research working proof}
\date{17 November 2026}
\begin{document}
\maketitle

For eight retained knots of meridional rank four, every reflective
augmentation has rank at most three. All eight have explicit reflective augmentations of rank exactly
three, in prime cyclotomic fields. They also admit explicit
rank-three Coxeter triangle quotients, so their maximum Coxeter
reflection-quotient rank is exactly three. Every
$SL_2\C$ character of each double branched cover is invariant under
the deck involution. We give direct linear-algebra proofs of the
upper bound and this character consequence.

Reflective augmentations are a specific commutative unoriented cord
quotient of knot contact homology. No upper bound on unrestricted
augmentation rank is claimed. The character consequence is a concrete
application, with an independent proof, of the conclusion of
Cornwell's conditional theorem \cite[Theorems 1.1 and 5.5]{cornwell}.
The uniform arguments and exact internal computations below have not
received formal verification or external specialist review. No global
priority claim is made.

\section{The exact cord-matrix convention}

We use Cornwell's positive sign convention for the unoriented abelian
cord ring $HC_0^{\mathrm{ab}}(K)$ \cite[Definitions 3.1, 3.3 and 3.4]{cornwell}.
For a diagram with arcs $0,\ldots,n-1$, an evaluation is specified by a
symmetric matrix $C=(c_{ij})$ with
\begin{equation}\label{eq:cord}
 c_{ii}=2,\qquad c_{\ell j}+c_{\ell k}=c_{\ell i}c_{ij}
 \quad\text{for every crossing }(i,j,k)\text{ and every }\ell.
\end{equation}
Here $i$ is the over-arc and $j,k$ are the two under-arcs.
The diagrammatic presentation theorem quoted in
\cite[Theorem 3.2]{cornwell}, with credit to Ng, identifies these
solutions with evaluations of that cord ring.
A reflective augmentation is an augmentation factoring through this
quotient, and its rank is $\rank C$. We import that exact presentation
and terminology; the upper-bound argument below is proved directly
from \eqref{eq:cord}, without an augmentation--representation theorem.
Write $a_{\mathrm{ref}}(K)$ for the largest reflective augmentation rank.

Let
\[
\mathcal K_8=\{13a_{373},13a_{375},13a_{391},13a_{392},
13a_{3264},13a_{3266},13a_{3275},13a_{3276}\}.
\]
The literal diagrams, arc numbering and original four-seed systems are
those of the retained seven-sheet certificates and universal reflection
obstruction. Knot identity and planarity are established by those
separate certificates, rather than inferred from the labels here.

\section{A Gram quotient converts cord matrices into reflections}

\begin{lemma}\label{lem:gram}
Let $k$ be any field of characteristic different from two, and let $C$
be a symmetric cord matrix satisfying \eqref{eq:cord} on one of the
retained diagrams. On
\[
                 V=k^n/\ker C,\qquad
                 B([x],[y])=x^{\mathsf T}Cy,
\]
the classes $v_i=[e_i]$ have norm two and satisfy
\[
                   v_j+v_k=c_{ij}v_i,\qquad
                   R_{v_i}(v_j)=-v_k,
                   \quad R_u(x)=x-B(u,x)u.
\]
The reflections $R_{v_i}$ give a knot-group representation. The
original four seed classes span $V$, whose dimension is $\rank C$.
\end{lemma}
\begin{proof}
The displayed form is well defined and nondegenerate on the quotient.
It has $B(v_i,v_i)=2$, so no $v_i$ vanishes. For each crossing,
\eqref{eq:cord} says
\[
 B(v_\ell,v_j+v_k-c_{ij}v_i)=0\qquad\text{for every }\ell.
\]
The vector therefore lies in $\ker C$ before taking the quotient,
and is zero afterwards. The reflection formula follows immediately.
Thus
\[
 R_{v_i}R_{v_j}R_{v_i}=R_{R_{v_i}v_j}=R_{-v_k}=R_{v_k},
\]
which is the involutory form of either sign of a Wirtinger relation.
Every retained propagation step expresses the new vector as a linear
combination of two known vectors. Hence the four original seed
classes span every $v_i$, and these classes span $V$.
\end{proof}

\begin{theorem}\label{thm:upper}
For every $K\in\mathcal K_8$, every symmetric cord matrix over a
field of characteristic different from two has rank at most three.
In particular,
\[
                    a_{\mathrm{ref}}(K)\le3<4=\mr(K)=b(K).
\]
\end{theorem}
\begin{proof}
Lemma~\ref{lem:gram} gives a common-form reflection representation
with the original four seed classes spanning a space of dimension
$\rank C$. The retained universal reflection theorem says their normal
directions are dependent for each of these eight diagrams, even if
an intermediate Gram form is singular. Thus $\rank C\le3$.
The meridional-rank and bridge equalities are supplied independently
by the retained seven-sheet cover and four-seed upper certificates.
The explicit finite-group epimorphisms in
\path{TERNARY_AFFINE_QUOTIENTS.tex} provide a second proof of these
equalities, without using a cover-group rank certificate.
\end{proof}

This proof never inverts a cord matrix or assumes that it is a Gram
matrix of independent vectors before passing to its radical quotient.
It concerns every complex solution of the cord equations, not just
real, positive definite, or finite-image solutions.

\section{Exact positive certificates and the remaining lower bound}

\begin{theorem}\label{thm:positive}
For every $K\in\mathcal K_8$,
\[
                      a_{\mathrm{ref}}(K)=3<4=\mr(K)=b(K).
\]
\end{theorem}
\begin{proof}
The following construction gives exact rank-three matrices. In each
row of the table, use either knot name in the left column and let
$p$ and the eight colours be as displayed. Columns are the reduced
arc labels obtained by identifying $4$ with $0$ and $7$ with $6$.
\[
\begin{array}{c|r|rrrrrrrr}
 K&p&0&5&6&8&9&10&11&12\\\hline
13a_{373},13a_{3264}&7&0&2&1&1&3&6&1&5\\
13a_{375},13a_{3266}&37&0&30&1&6&17&22&15&34\\
13a_{391},13a_{3275}&37&0&8&1&4&23&16&21&32\\
13a_{392},13a_{3276}&7&0&6&1&3&0&2&5&0
\end{array}
\]
These satisfy $f_j+f_k=2f_i\pmod p$ at each of the eight reduced
crossings. In particular, $f_0=0$ and $f_6=1$.

Let $\zeta$ be a primitive $p$th root of unity and put
$x=\zeta+\zeta^{-1}$. On a free three-dimensional space with basis
$E_0,E_1,V$, take the symmetric form
\[
 G=\begin{pmatrix}2&1&x\\1&2&0\\x&0&2\end{pmatrix}
 \quad\text{for }13a_{373},13a_{375},13a_{391},13a_{392},
\]
and
\[
 G=\begin{pmatrix}2&1&0\\1&2&-x\\0&-x&2\end{pmatrix}
 \quad\text{for }13a_{3264},13a_{3266},13a_{3275},13a_{3276}.
\]
Put $U=E_0-E_1$. In both cases $U,V$ have norm two and
$B(U,V)=x$. Define planar roots by
\[
                  r_0=U,\quad r_1=V,\quad
                  r_{k+1}=xr_k-r_{k-1}.
\]
The characteristic roots of this recurrence are $\zeta,\zeta^{-1}$,
so $r_{k+p}=r_k$. Their norms are two and
$R_{r_i}(r_j)=-r_{2i-j}$.
For arcs outside $\{1,2,3\}$ assign the raw vector $r_{f_a}$,
using the two identifications above. Assign raw vectors $E_0$ to
arc $1$, $E_1$ to arc $3$, and $E_1$ or $E_0$ to arc $2$ in the
first or second Gram convention, respectively.
Multiply these thirteen raw vectors by the signs
\[
                     (1,-1,1,-1,-1,-1,-1,1,1,1,1,1,1).
\]
Denote the resulting vectors by $v_a$ and set $c_{ab}=B(v_a,v_b)$.
The eight reduced Fox equations, the local norm-two $A_2$ relations,
and the displayed signs give $v_k=-R_{v_i}v_j$ at every crossing.
The checker reconstructs all vectors and verifies every crossing-row
identity in \eqref{eq:cord}, rather than assuming this sign compatibility.

The factorization $C=RGR^{\mathsf T}$ gives rank at most three.
The principal minor at arcs $(1,2,6)$ in the first convention or
$(1,3,6)$ in the second is
\[
                            \det G=2(3-x^2)\ne0.
\]
For nonvanishing, $x^2=3$ would imply
$\zeta^4-\zeta^2+1=0$, and hence $\zeta^{12}=1$, impossible for
primitive order $7$ or $37$. Thus the matrices have rank exactly
three over $\Q(\zeta)$, and any complex embedding gives the required
reflective augmentation. Theorem~\ref{thm:upper} proves maximality.
\end{proof}

Four additional independent controls are retained over
$\Q(\phi)$, where $\phi^2=\phi+1$, for
$13a_{373},13a_{392},13a_{3264},13a_{3276}$.
Each gives thirteen explicit vectors in $\Q(\phi)^4$ with norm four,
sets $c_{ij}=(r_i\cdot r_j)/2$, and verifies all cord equations and
rank three by exact field elimination. These witnesses were found
from reflection colourings followed by a separate signed cord check.
They are not needed for the uniform cyclotomic construction.

\section{Explicit rank-three Coxeter triangle quotients}

\begin{theorem}\label{thm:triangle}
Every $K\in\mathcal K_8$ has a meridian-preserving epimorphism onto
the Coxeter triangle group $W(2,3,p)$, where $p=7$ or $37$ is the
prime assigned in the table above. Together with the universal
reflection obstruction, this shows that its maximum Coxeter
reflection-quotient rank is exactly three.
\end{theorem}
\begin{proof}
Use generators $A,B,C$ of order two. In the first Gram convention,
let the target presentation be
\[
       (AB)^3=(BC)^2=(AC)^p=1,
\]
and in the second let it be
\[
       (AB)^3=(AC)^2=(BC)^p=1.
\]
These are two labellings of the same rank-three Coxeter triangle
system. Put $S=BAB=ABA$, the equality following from $(AB)^3=1$.
In either presentation, $(CS)^p=1$: in the first, $B$ commutes with
$C$ and conjugates $CA$ to $CS$; in the second, $A$ commutes with
$C$ and conjugates $CB$ to $CS$.

Set $q=CS$. The reflections $t_i=q^iS$, indexed modulo $p$, satisfy
\[
                           t_i t_j t_i=t_{2i-j}.
\]
Assign $t_{f_a}$ to every reduced arc in the table; the eight Fox
congruences give all reduced Wirtinger relations. On the five local
crossings, assign $A$ to arc $1$ and $B$ to arc $3$. Assign $B$ to
arc $2$ in the first convention and $A$ in the second, while arcs
$0,4$ are $S$ and arcs $6,7$ are $C$. The relations of the $A_2$
pair $A,B$ and the indicated commuting pair verify those remaining
crossings. Every assigned image is a Coxeter reflection: $S$ is a
conjugate of a simple generator, and since $p$ is odd all $q^iS$
are conjugate to $S$ in the dihedral subgroup.

The images include all three standard generators $A,B,C$, so the
homomorphism is onto the abstract Coxeter group. This argument does
not infer an abstract quotient merely from a matrix image, and uses
no faithfulness or specialization assertion. The universal
reflection theorem gives the matching upper bound three on Coxeter
reflection-quotient rank.
\end{proof}

The groups $W(2,3,7)$ and $W(2,3,37)$ are hyperbolic triangle
Coxeter groups. These quotients are consistent with the absence of
sharp finite Coxeter targets: their rank-three lower bounds still
fall short of the seven-sheet lower bound four.

\section{Deck invariance of every double-cover character}

\begin{theorem}\label{thm:deck}
For every $K\in\mathcal K_8$ and every
$\rho:\pi_1(\Sigma_2K)\to SL_2\C$, the characters of $\rho$ and
$\rho\circ\tau_*$ agree, where $\tau$ is the deck involution.
Equivalently,
\[
                          X^\tau(\Sigma_2K)=X(\Sigma_2K).
\]
\end{theorem}
\begin{proof}
In the meridian-square orbifold group, fix a meridian $m_0$ and put
$a_i=m_i m_0$, with $a_0=1$. Reidemeister--Schreier rewriting gives
the double-cover presentation
\[
                 a_k=a_i a_j^{-1}a_i\quad\text{at }(i,j,k),
                  \qquad \tau_*(a_i)=a_i^{-1}.
\]
The second transversal rewrite is the inverse of the first, so no
additional relation is being omitted. Set $A_i=\rho(a_i)$ and
consider $M_2(\C)$ with the nondegenerate symmetric form
\[
                          B(X,Y)=\tr(X\adj Y).
\]
It has $B(A_i,A_i)=2$. The generic two-by-two identity
\begin{equation}\label{eq:adj}
 X\adj(Y)X=B(X,Y)X-\det(X)Y
\end{equation}
shows
\[
                   A_k=B(A_i,A_j)A_i-A_j=-R_{A_i}(A_j).
\]
As in Lemma~\ref{lem:gram}, these normals define a knot-group
reflection representation, and the four seed matrices span all
$A_i$. The universal reflection theorem therefore gives
\[
                         E=\langle A_i\rangle,\qquad\dim E\le3.
\]
Since $I=A_0\in E$, all traceless parts
$X_i=A_i-\frac12\tr(A_i)I$ lie in a subspace
$W\subset\mathfrak{sl}_2(\C)$ of dimension at most two.

Suppose first that $\rho$ is irreducible. If $\dim W\le1$, all
$A_i$ are polynomials in a common matrix and have a common eigenline,
a contradiction. Thus $\dim W=2$.
The trace form $\tr(XY)$ is nondegenerate on $W$. Otherwise a nonzero
radical vector $N\in W$ satisfies $\tr(N^2)=0$. As $N$ is traceless,
Cayley--Hamilton gives $\det N=0$, so $N$ is nonzero nilpotent.
Conjugate it to an upper triangular matrix supported at entry $(1,2)$.
The equalities $\tr(NX_i)=0$ force every $X_i$ to be upper triangular,
contradicting irreducibility.

The orthogonal complement of $W$ in the nondegenerate
three-dimensional trace space $\mathfrak{sl}_2(\C)$ is therefore a
nonisotropic line. Choose a nonzero matrix $Z$ on this line. It is
invertible, and the traceless two-by-two identity
\[
                         ZX+XZ=\tr(ZX)I
\]
gives $ZX_i=-X_iZ$. Consequently
\[
                 ZA_iZ^{-1}=\tfrac12\tr(A_i)I-X_i
                              =\adj A_i=A_i^{-1}.
\]
Rescale $Z$ to determinant one, which is possible over $\C$.
The resulting matrix conjugates $\rho$ to $\rho\circ\tau_*$ on the
generators and therefore on the full group.

If $\rho$ is reducible, replace it by its diagonal semisimplification,
which has the same character. Its diagonal characters are $\eta$
and $\eta^{-1}$. Since $\tau_*$ inverts every $a_i$, it interchanges
these two diagonal characters. Swapping the two diagonal entries is
conjugation in $SL_2\C$, so the character is again unchanged.
\end{proof}

Cornwell's Theorem 5.5 proves this conclusion under the general
hypothesis $a_{\mathrm{ref}}(K)\le3$. The proof above independently
derives its first character conclusion on the eight retained knots.
No new general equivalence with arbitrary character varieties is
claimed, and no assertion about their reduced coordinate rings is
needed here.

\section{Verification and scope}

The standard-library checker reconstructs the canonical arc components
of all eight literal diagrams. It verifies 2,028 exact crossing-row identities across eight prime
cyclotomic witnesses and four independent quadratic-field controls,
exact ranks or nonzero principal minors, and the cyclotomic root
recurrences. The reduced colours and local $A_2$ pattern are checked
against the literal diagrams. Dihedral conjugation is separately
checked on all pairs of polygon reflections at orders seven and
thirty-seven.
It separately verifies \eqref{eq:adj}, symmetry and nondegeneracy of
the matrix form, and the traceless anticommutator identity as generic
integer polynomial identities. A deliberately corrupted symmetric
cord matrix is rejected.

The uniform upper bound uses the written Gram quotient and the
retained universal reflection proof. The deck statement also uses
the written presentation and irreducible/reducible linear-algebra
arguments. The finite controls do not formally verify those arguments.
Unrestricted augmentation rank, external review and complete
worldwide priority remain open.

\begin{thebibliography}{9}
\bibitem{cornwell} C. R. Cornwell,
Character varieties of knot complements and branched double-covers
via the cord ring, arXiv:1509.04962v1 (2015), Sections 3--5,
Theorem 3.2, Definitions 3.3--3.4, and Theorems 1.1 and 5.5.
\url{https://arxiv.org/abs/1509.04962v1}.
\end{thebibliography}
\end{document}
